\begin{abstract}
The viscoelastic gel-point of a polymer network is commonly treated as a fixed property of chemistry and concentration. Here we demonstrate, through three-dimensional random walk simulations on site-percolation lattices, that the gel-point is instead controlled by the growth dynamics of the forming network and can be continuously programmed via a single nucleation-density parameter. Probe-particle diffusion is modelled by random walkers confined to the unoccupied (void) sublattice of an $L \times L \times L$ cubic lattice ($L = 500$); the gel-point $p_c'$ is identified by the Winter--Chambon criterion $\alpha = 0.5$, where $\langle r^2(t) \rangle \propto t^\alpha$. For random (Bernoulli) percolation networks we find $p_c' \approx 0.683 \approx 1 - p_c$, quantitatively identifying the rheological gel-point with the loss of spanning connectivity in the void phase. A single nucleation-density parameter $\kappa \in [0,1]$ — the fraction of newly occupied sites grown from the cluster frontier rather than placed as independent nuclei — then shifts this gel-point continuously and monotonically over $\Delta p_c' \approx 0.31$ ($p_c' \in [0.683, 0.993]$), from the random limit ($\kappa = 0$) towards a compact single-seed limit ($\kappa \to 1$, $p_c' \to 1$). As growth becomes more compact the transition broadens (up to $\approx 3.4$ times the random-limit width) and probe particles remain persistently sub-diffusive ($\alpha \approx 0.3$--$0.4$) above $p_c'$; a many-nuclei, adjacency-templated morphology sits at $p_c' \approx 0.886$ as one point along this continuum rather than a distinct gelation class. Frequency-domain viscoelastic spectra $G'(\omega)$ and $G''(\omega)$ derived via the Generalised Stokes--Einstein Relation confirm the $\delta(\omega) = 45^{\circ}$ Winter--Chambon signature at $p_c'$ across these morphologies. These results establish nucleation density as a structural design parameter for programming the gel-point in network-forming materials.
\end{abstract}
